\section{Session Dynamics and Persistence Mechanism}
\label{sec:session-persistence}

\Cref{sec:formal-model} fixed the state object, policies, and closure map; here those pieces are composed into an end-to-end session update: retrieve a slice, expand local open structure from each agent, merge into a candidate set, close under $\Pi$.
Unless noted, operators are partial relations; determinism is not assumed.
The \emph{reference merge policy} $\PsiMerge^{\mathrm{ref}}$ (\Cref{def:ref-merge}) completes the triple $(\PsiConf^{\mathrm{ref}},\PsiMerge^{\mathrm{ref}},\Pi_{\mathrm{prov}}^{\mathrm{ref}})$ fixed in \Cref{sec:reference-policies}.

\subsection{Retrieval / slice operator}
\label{sec:op-slice}

\begin{definition}[Slice operator]
\label{def:slice-op}
$\SliceOp_t : Q \times \{\GlobalState_t\} \to \{\SessionState_t(q)\}$ produces a session slice satisfying \Cref{def:session-slice}.
The selection of $V^{\mathrm{sel}}_t$ is read-only copy from $\GlobalState_t$.
Intuitively, slicing is ``load context'': it materializes a working subgraph without yet committing new persistent thoughts.
\end{definition}

\subsection{Stem and expansion}
\label{sec:op-expand}

\begin{definition}[Stem map]
\label{def:stem}
$\StemOp_t : A \times \SessionState_t(q) \to V_t \cup \mathcal{P}_{\mathrm{fin}}(V_t)$ as before.
Intuitively, a stem is where an agent attaches new reasoning---a hook vertex or small hook set in the slice.
\end{definition}

\begin{definition}[Local session expansion]
\label{def:expand-op}
$\ExpandOp_t$ maps $(a,\SessionState_t(q),s^{(a)}_t)$ to a finite open fragment $(\Delta V^{(a)}_t,\Delta \mathcal{E}^{(a)}_t)$ with session scope.
Intuitively, expansion is local hypothesis formation: new open vertices and edges that are not yet globally persistent.
\end{definition}

\subsection{Merge operator and merge policy}
\label{sec:op-merge}

\begin{definition}[Merge policy]
\label{def:merge-policy}
A \emph{merge policy} $\PsiMerge$ specifies:
\begin{itemize}[leftmargin=*]
  \item an \emph{equivalence} relation $\sim$ on proposed vertex identifiers (duplicate detection / canonicalization);
  \item whether merge is \emph{order-dependent} (tuple of agents ordered) or \emph{commutative} up to $\sim$;
  \item how to combine parallel edges (union, multiset union, or quotient).
\end{itemize}
Intuitively, $\PsiMerge$ specifies how independent agents' open fragments are aligned, deduplicated, and welded into one candidate graph.
\end{definition}

\begin{definition}[Merge operator]
\label{def:merge-op}
Given fragments $(\Delta V^{(a)}_t,\Delta \mathcal{E}^{(a)}_t)_{a\in A}$, a \emph{merge} operator $\MergeOp_t^{\PsiMerge}$ produces a single candidate set $\CandSet_t=(V^{\mathrm{c}}_t,E^{\mathrm{c}}_t)$ by:
\begin{enumerate}[label=\arabic*.,leftmargin=*]
  \item forming the disjoint union of fragments;
  \item quotienting vertices by $\sim$ (if $\PsiMerge$ identifies duplicates);
  \item taking the union of edge sets, relabeling endpoints under the quotient;
  \item optionally inserting $\mathrm{conf}$ edges between mutually inconsistent survivors when $\PsiMerge$ declares incompatibility.
\end{enumerate}
For fixed $\PsiMerge$, $\MergeOp_t^{\PsiMerge}$ may still be \emph{set-valued} if multiple maximal conflict resolutions are permitted; a complete system description then either applies $\CloseOp_t^{\Pi}$ to each feasible candidate or specifies a selection rule that chooses one candidate set.
Intuitively, merge reconciles parallel work into a single batch of candidates; closure then decides global acceptance.
\end{definition}

\begin{definition}[Reference merge policy]
\label{def:ref-merge}
The \emph{reference merge policy} $\PsiMerge^{\mathrm{ref}}$ specifies:
\begin{enumerate}[label=(\roman*),leftmargin=*]
  \item an equivalence relation $\sim$ on vertex identifiers, fixed per session, that canonicalizes duplicate hypotheses (e.g.\ same synthetic label after hash-based renaming);
  \item \textbf{commutative} merge up to $\sim$ (no agent order dependence);
  \item disjoint union of fragments, followed by quotient of vertices by $\sim$;
  \item edge set as the union of images of edges under the quotient map on endpoints;
  \item \textbf{retain tension:} if two $\sim$-classes contain vertices $v,w$ with $(v,w)\in\ConfRel_t$ after union, insert a single $\mathrm{conf}$ edge between chosen representatives (unless already present), so conflicting candidates survive for $\PsiConf$ rather than being silently dropped.
\end{enumerate}
Under $\PsiMerge^{\mathrm{ref}}$, $\MergeOp_t^{\PsiMerge^{\mathrm{ref}}}$ is single-valued on finite inputs once $\sim$ is fixed.
\end{definition}

\textbf{Merge before closure.} The pipeline used in this paper is: expand $\to$ merge to $\CandSet_t$ $\to$ $\CloseOp_t^{\Pi}(\GlobalState_t,\CandSet_t)$.
Merge does not write to $\GlobalState_t$; only $\CloseOp_t^{\Pi}$ does.

\subsection{Conflict policy interaction}
\label{sec:op-conflict}

$\PsiConf$ (\Cref{def:conflict-policy}) acts on $(\ThoughtGraph_t,\CandSet_t,\ConfRel_t)$ after candidate attachment to the ambient graph is specified for the purpose of evaluating $\ConfRel_t$ on candidate vertices.

\subsection{Closure reference}
\label{sec:op-close-repeat}

$\CloseOp_t^{\Pi}$ is \Cref{def:closure-op}.
Provenance policy $\Pi_{\mathrm{prov}}$ decides which stem/provenance edges become part of $\mathcal{E}^{\mathrm{g}}_{t+1}$ versus ephemeral audit logs (external to $\GlobalState$).

\subsection{End-to-end composition}
\label{sec:compose}

\begin{align}
  \SessionState_t(q) &= \SliceOp_t(q,\GlobalState_t), \\
  \CandSet_t &\in \MergeOp_t^{\PsiMerge}\Bigl(\bigl(\ExpandOp_t(a,\SessionState_t(q),\StemOp_t(a,\SessionState_t(q)))\bigr)_{a\in A}\Bigr), \\
  (\GlobalState_{t+1},\mathsf{status}) &= \CloseOp_t^{\Pi}(\GlobalState_t,\CandSet_t).
\end{align}
Ingest may prepend: $\IngestOp_t$ updates $\Vlayer{0}$ before $\SliceOp_t$.

\begin{figure}[t]
  \centering
  % Figure 3 — Session lifecycle pipeline (PROJECT_SPEC.md 11A)
% Dashed group uses `fit` around agent boxes so the frame does not cross interior text.
\resizebox{\linewidth}{!}{%
\begin{tikzpicture}[
  font=\small,
  box/.style={draw, rounded corners, minimum height=0.9cm, minimum width=2.0cm, align=center},
  arr/.style={-{Latex[length=2mm]}, thick},
]

\node[box] (q) at (0,0) {\textbf{Query}};
\node[box] (sl) at (2.6,0) {\textbf{Slice}\\$\SliceOp$};
\node[box] (a1) at (5.3,0.75) {\textbf{Agent 1}\\stem};
\node[box] (a2) at (5.3,-0.75) {\textbf{Agent 2}\\stem};
\node[box] (m) at (8.1,0) {\textbf{Debate/Merge}\\$\MergeOp$};
\node[box] (cl) at (10.9,0) {\textbf{Close}\\$\CloseOp^{\Pi}$};
\node[box] (rw) at (13.4,0) {\textbf{Persist}\\$\GlobalState_{t+1}$};

\draw[arr] (q) -- (sl);
\draw[arr] (sl) |- (a1);
\draw[arr] (sl) |- (a2);
\draw[arr] (a1) -| node[near start, above, xshift=-6pt] {open} (m);
\draw[arr] (a2) -| (m);
\draw[arr] (m) -- (cl);
\draw[arr] (cl) -- (rw);

\node[draw, dashed, rounded corners, inner sep=9pt, fit=(a1)(a2)] (fitagents) {};
\node[below=5pt of fitagents.south, font=\footnotesize] {parallel expansion (schematic)};

\end{tikzpicture}%
}

  \caption{Session pipeline: slice, parallel expand, merge into $\CandSet_t$ (e.g.\ $\PsiMerge^{\mathrm{ref}}$, \Cref{def:ref-merge}), then $\CloseOp_t^{\Pi}$ (\Cref{def:closure-op}). Only closure updates $\GlobalState_{t+1}$.}
  \label{fig:pipeline}
\end{figure}

\begin{algorithm}[t]
  \caption{Session update at $t$ for query $q$ (policy-relative)}
  \label{alg:main}
  \begin{algorithmic}[1]
    \State $\SessionState_t(q) \gets \SliceOp_t(q,\GlobalState_t)$
    \For{each $a\in A$}
      \State $(\Delta V^{(a)}_t,\Delta \mathcal{E}^{(a)}_t) \gets \ExpandOp_t(a,\SessionState_t(q),\StemOp_t(a,\SessionState_t(q)))$
    \EndFor
    \State choose $\CandSet_t \in \MergeOp_t^{\PsiMerge}\bigl((\Delta V^{(a)}_t,\Delta \mathcal{E}^{(a)}_t)_{a\in A}\bigr)$ if set-valued
    \State $(\GlobalState_{t+1},\mathsf{status}) \gets \CloseOp_t^{\Pi}(\GlobalState_t,\CandSet_t)$
  \end{algorithmic}
\end{algorithm}

\subsection{Design-level failure modes}
\label{sec:failure-modes}

If $\MergeOp_t^{\PsiMerge}$ is set-valued and no selection rule is fixed, the pipeline need not advance---a form of \emph{merge deadlock} at the design level.
Other obstruction categories: unsupported candidates after $\SigmaAdm$, conflict pairs after $\PsiConf$, or empty $\CandSet^{\mathrm{pass}}$.
